\documentclass[12pt,a4paper]{article}

\usepackage[T1]{fontenc}
\usepackage[utf8]{inputenc}
\usepackage[romanian]{babel}
\usepackage{geometry}
\usepackage{amsmath}

\geometry{margin=2.5cm}

\title{Tutoriat 3\\Cautare Adversariala - Algoritmul MIN-MAX}
\author{}
\date{}

\begin{document}

\maketitle


\section{Algoritmul MiniMax}

Algoritmul MiniMax este una dintre cele mai importante metode de căutare adversarială utilizate în Inteligența Artificială.

Scopul algoritmului este determinarea mutării optime pentru un jucător presupunând că adversarul joacă întotdeauna perfect.

MiniMax este utilizat în special pentru jocuri:

\begin{itemize}
\item cu doi jucători;
\item cu informație completă;
\item deterministe;
\item cu sumă zero.
\end{itemize}

Exemple:

\begin{itemize}
\item șah;
\item dame;
\item X și 0;
\item Connect Four;
\item Reversi.
\end{itemize}

\subsection{Ideea de bază}

Algoritmul presupune existența a doi jucători:

\begin{itemize}
\item MAX - încearcă să maximizeze scorul;
\item MIN - încearcă să minimizeze scorul.
\end{itemize}

Fiecare jucător alege întotdeauna mutarea cea mai avantajoasă pentru propriul interes.

MAX presupune că MIN va juca perfect.

MIN presupune că MAX va juca perfect.

\subsection{Principiul de funcționare}

MiniMax explorează arborele de joc până la stările terminale.

Valorile stărilor terminale sunt apoi propagate către rădăcină.

Regula este:

\begin{itemize}
\item într-un nod MAX se alege valoarea maximă;
\item într-un nod MIN se alege valoarea minimă.
\end{itemize}

Astfel se poate determina cea mai bună mutare disponibilă pentru MAX.

\subsection{Exemplu}

Considerăm următorul arbore:

\begin{verbatim}
                               MAX
                    /-------------|-------------\
                  MIN            MIN           MIN
                /-----\        /-----\       /-----\
              MAX    MAX     MAX    MAX    MAX    MAX
             /  \    /  \   /  \   /  \   /  \   /  \
            3    5  6    9 1    2 0    7 4    8 2    6
\end{verbatim}


Pasul 1: evaluăm nodurile MAX.

\[
max(3,5)=5
\]

\[
max(6,9)=9
\]

\[
max(1,2)=2
\]

\[
max(0,7)=7
\]

\[
max(4,8)=8
\]

\[
max(2,6)=6
\]

Obținem:

\begin{verbatim}
                MAX
           /------|------\
         MIN     MIN     MIN
        /  \    /  \    /  \
       5    9  2    7  8    6
\end{verbatim}

Pasul 2: evaluăm nodurile MIN.

\[
min(5,9)=5
\]

\[
min(2,7)=2
\]

\[
min(8,6)=6
\]

Obținem:

\begin{verbatim}
               MAX
            /   |   \
           5    2    6
\end{verbatim}

Pasul 3: evaluăm rădăcina.

\[
max(5,2,6)=6
\]

Valoarea finală a rădăcinii este:

\[
6
\]

Prin urmare, MAX va selecta ramura care conduce la valoarea 6.

\subsection{Intuiție}

MAX încearcă să obțină cel mai mare scor posibil.

MIN încearcă să împiedice acest lucru.

Rezultatul final reprezintă cel mai bun scor pe care MAX îl poate garanta indiferent de răspunsul adversarului.

\subsection{Strategia optimă}

MiniMax nu caută mutarea care oferă cel mai bun rezultat posibil.

MiniMax caută mutarea care oferă cel mai bun rezultat garantat.

Aceasta este o diferență foarte importantă.

Algoritmul presupune întotdeauna că adversarul va alege cea mai nefavorabilă mutare pentru noi.

\subsection{Avantaje}

\begin{itemize}
\item produce decizii optime;
\item garantează cea mai bună strategie în jocurile cu sumă zero;
\item este simplu de înțeles și implementat;
\item stă la baza multor algoritmi moderni pentru jocuri.
\end{itemize}


\subsection{Dezavantaje}

\begin{itemize}
\item necesită explorarea unui număr foarte mare de noduri;
\item consumă mult timp pentru arbori mari;
\item consumă memorie pentru jocuri complexe;
\item devine impracticabil pentru jocuri cu factori mari de ramificare.
\end{itemize}

\subsection{Concluzie}

MiniMax reprezintă algoritmul fundamental pentru căutarea adversarială.

Ideea sa principală este propagarea valorilor din frunzele arborelui către rădăcină folosind operatorii maxim și minim.

Din cauza costului ridicat de calcul, în practică MiniMax este adesea combinat cu funcții de evaluare și cu tehnica Alpha-Beta Pruning.

























\section{Formalizarea algoritmului MiniMax}

Pentru a putea implementa algoritmul MiniMax, este necesară descrierea sa într-o formă matematică riguroasă.

Formalizarea definește modul în care se calculează valoarea fiecărui nod din arborele de joc.

\subsection{Ipoteze}

Algoritmul MiniMax presupune că:

\begin{itemize}
\item există doi jucători: MAX și MIN;
\item jocul este cu sumă zero;
\item ambii jucători joacă optim;
\item jocul este determinist;
\item informația despre stare este completă.
\end{itemize}

\subsection{Valoarea unei stări}

Fie o stare $s$.

Vom nota prin:

\[
V(s)
\]

valoarea asociată stării respective.

Această valoare reprezintă rezultatul pe care MAX îl poate garanta dacă ambii jucători joacă optim.

\subsection{Stări terminale}

Dacă starea este terminală, valoarea ei este dată direct de funcția de utilitate.

\[
V(s)=U(s)
\]

unde:

\begin{itemize}
\item $U(s)$ reprezintă funcția de utilitate;
\item $s$ este o stare terminală.
\end{itemize}

Exemplu:

\[
U(s)=1
\]

pentru victoria lui MAX.

\[
U(s)=0
\]

pentru remiză.

\[
U(s)=-1
\]

pentru victoria lui MIN.

\subsection{Noduri MAX}

Dacă este rândul lui MAX să mute, acesta va alege succesorul cu valoarea cea mai mare.

Formal:

\[
V(s)=\max_{s' \in Succ(s)} V(s')
\]

unde:

\begin{itemize}
\item $Succ(s)$ reprezintă mulțimea succesorilor lui $s$;
\item $s'$ reprezintă un succesor al lui $s$.
\end{itemize}

Exemplu:

\begin{verbatim}
      MAX
     / | \
    2  7  4
\end{verbatim}

Atunci:

\[
V(MAX)=\max(2,7,4)=7
\]

\subsection{Noduri MIN}

Dacă este rândul lui MIN să mute, acesta va alege succesorul cu valoarea cea mai mică.

Formal:

\[
V(s)=\min_{s' \in Succ(s)} V(s')
\]

Exemplu:

\begin{verbatim}
      MIN
     / | \
    2  7  4
\end{verbatim}

Atunci:

\[
V(MIN)=\min(2,7,4)=2
\]

\subsection{Definiția recursivă a MiniMax}

Algoritmul poate fi definit recursiv astfel:

\[
MiniMax(s)=
\begin{cases}
U(s), & \text{dacă } s \text{ este terminală} \\\\
\max\limits_{s' \in Succ(s)} MiniMax(s'),
& \text{dacă } s \text{ este nod MAX} \\\\
\min\limits_{s' \in Succ(s)} MiniMax(s'),
& \text{dacă } s \text{ este nod MIN}
\end{cases}
\]

Această definiție constituie baza tuturor implementărilor MiniMax.

\subsection{Calculul valorilor}

Calculul începe întotdeauna din frunzele arborelui.

Pașii sunt:

\begin{enumerate}
\item evaluarea stărilor terminale;
\item calculul valorilor nodurilor MIN;
\item calculul valorilor nodurilor MAX;
\item propagarea valorilor până la rădăcină.
\end{enumerate}

Procesul este numit propagare inversă (back-propagation of values).

\subsection{Exemplu}

Considerăm arborele:

\begin{verbatim}
               MAX
            /       \
          MIN       MIN
         /  \       /  \
        3    5     2    9
\end{verbatim}

Calculăm:

\[
V(MIN_1)=\min(3,5)=3
\]

\[
V(MIN_2)=\min(2,9)=2
\]

Apoi:

\[
V(MAX)=\max(3,2)=3
\]

Prin urmare:

\[
MiniMax(MAX)=3
\]

\subsection{Semnificația valorii rădăcinii}

Valoarea rădăcinii reprezintă rezultatul pe care MAX îl poate garanta indiferent de modul în care adversarul încearcă să îl împiedice.

Aceasta nu este neapărat cea mai bună valoare posibilă, ci cea mai bună valoare garantată.

\subsection{Concluzie}

Formalizarea MiniMax transformă problema jocului într-o problemă de optimizare recursivă.

MAX aplică operația de maximizare, MIN aplică operația de minimizare, iar valorile sunt propagate din frunzele arborelui către rădăcină până la obținerea strategiei optime.











\section{Implementarea algoritmului MiniMax}

După definirea formală a algoritmului MiniMax, următorul pas este implementarea acestuia.

MiniMax este un algoritm recursiv care explorează arborele de joc și propagă valorile de la frunze către rădăcină.

Ideea principală este:

\begin{itemize}
\item dacă nodul aparține lui MAX, se alege valoarea maximă;
\item dacă nodul aparține lui MIN, se alege valoarea minimă;
\item dacă s-a ajuns într-o stare terminală, se returnează utilitatea acesteia.
\end{itemize}

\subsection{Pseudocod}

\begin{verbatim}
MiniMax(stare, jucator)

    daca stare este terminala
        returneaza utilitatea(stare)

    daca jucator este MAX

        ceaMaiBunaValoare <- -infinit

        pentru fiecare succesor
            ceaMaiBunaValoare <-
                max(ceaMaiBunaValoare,
                    MiniMax(succesor, MIN))

        returneaza ceaMaiBunaValoare

    altfel

        ceaMaiBunaValoare <- +infinit

        pentru fiecare succesor
            ceaMaiBunaValoare <-
                min(ceaMaiBunaValoare,
                    MiniMax(succesor, MAX))

        returneaza ceaMaiBunaValoare
\end{verbatim}

\subsection{Implementare în Python}

O implementare simplificată a algoritmului MiniMax poate fi scrisă astfel:

\begin{verbatim}
def minimax(node, maximizing_player):

    if node.is_terminal():
        return node.utility()

    if maximizing_player:

        best_value = float("-inf")

        for child in node.children:
            value = minimax(child, False)
            best_value = max(best_value, value)

        return best_value

    else:

        best_value = float("inf")

        for child in node.children:
            value = minimax(child, True)
            best_value = min(best_value, value)

        return best_value
\end{verbatim}

\subsection{Explicația implementării}

Funcția primește:

\begin{itemize}
\item nodul curent;
\item informația dacă este rândul lui MAX sau al lui MIN.
\end{itemize}

Dacă nodul este terminal, se returnează valoarea utilității.

Dacă este rândul lui MAX:

\begin{itemize}
\item se explorează toți succesorii;
\item se calculează recursiv valorile acestora;
\item se păstrează valoarea maximă.
\end{itemize}

Dacă este rândul lui MIN:

\begin{itemize}
\item se explorează toți succesorii;
\item se calculează recursiv valorile acestora;
\item se păstrează valoarea minimă.
\end{itemize}

\subsection{Exemplu de arbore}

Considerăm arborele:

\begin{verbatim}
               MAX
            /       \
          MIN       MIN
         /  \       /  \
        3    5     2    9
\end{verbatim}

Calcul:

\[
MIN(3,5)=3
\]

\[
MIN(2,9)=2
\]

\[
MAX(3,2)=3
\]

Algoritmul va returna valoarea:

\[
3
\]

și va selecta ramura din stânga.

\subsection{Complexitate}

Pentru un factor de ramificare b și o adâncime d:

\[
O(b^d)
\]

Numărul de stări analizate crește exponențial cu adâncimea arborelui.

Acesta este principalul motiv pentru care MiniMax devine dificil de utilizat în jocuri foarte complexe.

\subsection{Observație}

În practică, MiniMax este aproape întotdeauna combinat cu:

\begin{itemize}
\item limitarea adâncimii de căutare;
\item funcții de evaluare;
\item Alpha-Beta Pruning.
\end{itemize}

Aceste tehnici reduc semnificativ timpul necesar pentru luarea unei decizii.










\section{Exemple MiniMax}

În această secțiune vom analiza mai multe exemple de aplicare a algoritmului MiniMax pentru a înțelege modul în care sunt calculate valorile nodurilor și cum este aleasă mutarea optimă.

\subsection{Exemplul 1}

Considerăm următorul arbore:

\begin{verbatim}
               MAX
            /       \
          MIN       MIN
         /  \       /  \
        3    5     2    9
\end{verbatim}

Calculăm valorile nodurilor MIN:

\[
MIN(3,5)=3
\]

\[
MIN(2,9)=2
\]

Apoi calculăm valoarea rădăcinii:

\[
MAX(3,2)=3
\]

Rezultat:

\[
V(MAX)=3
\]

Mutarea optimă este ramura stângă.

\subsection{Exemplul 2}

Considerăm un arbore cu un nivel suplimentar:

\begin{verbatim}
                    MAX
                 /       \
               MIN       MIN
              /  \       /  \
            MAX  MAX   MAX  MAX
            / \  / \   / \  / \
           3  5 2  6  1  4 7  8
\end{verbatim}

Calculăm valorile nodurilor MAX:

\[
MAX(3,5)=5
\]

\[
MAX(2,6)=6
\]

\[
MAX(1,4)=4
\]

\[
MAX(7,8)=8
\]

Obținem:

\begin{verbatim}
                MAX
             /       \
           MIN       MIN
          /  \       /  \
         5    6     4    8
\end{verbatim}

Calculăm valorile nodurilor MIN:

\[
MIN(5,6)=5
\]

\[
MIN(4,8)=4
\]

Apoi:

\[
MAX(5,4)=5
\]

Rezultatul final este:

\[
V(MAX)=5
\]

\subsection{Exemplul 3}

Considerăm arborele:

\begin{verbatim}
                               MAX
                    /-------------|-------------\
                  MIN            MIN           MIN
                /-----\        /-----\       /-----\
              MAX    MAX     MAX    MAX    MAX    MAX
             /  \    /  \   /  \   /  \   /  \   /  \
            3    5  6    9 1    2 0    7 4    8 2    6
\end{verbatim}

Evaluăm nodurile MAX:

\[
max(3,5)=5
\]

\[
max(6,9)=9
\]

\[
max(1,2)=2
\]

\[
max(0,7)=7
\]

\[
max(4,8)=8
\]

\[
max(2,6)=6
\]

Rezultatul devine:

\begin{verbatim}
                MAX
           /------|------\
         MIN     MIN     MIN
        /  \    /  \    /  \
       5    9  2    7  8    6
\end{verbatim}

Calculăm nodurile MIN:

\[
min(5,9)=5
\]

\[
min(2,7)=2
\]

\[
min(8,6)=6
\]

Obținem:

\begin{verbatim}
               MAX
            /   |   \
           5    2    6
\end{verbatim}

Valoarea rădăcinii:

\[
max(5,2,6)=6
\]

Prin urmare:

\[
V(MAX)=6
\]

MAX va alege a treia ramură.

\subsection{Observație}

În toate exemplele calculul se realizează de jos în sus:

\begin{enumerate}
\item evaluăm frunzele;
\item calculăm nodurile MAX;
\item calculăm nodurile MIN;
\item determinăm valoarea rădăcinii.
\end{enumerate}

Această propagare a valorilor reprezintă mecanismul fundamental al algoritmului MiniMax.

\subsection{Concluzie}

MiniMax permite determinarea mutării optime presupunând că adversarul joacă perfect.

Valoarea calculată în rădăcina arborelui reprezintă cel mai bun rezultat pe care jucătorul MAX îl poate garanta indiferent de răspunsul adversarului.










\section{Proprietățile algoritmului MiniMax}

Algoritmul MiniMax permite determinarea mutării optime într-un joc cu doi jucători presupunând că ambii participanți joacă perfect.

Analiza proprietăților sale este importantă deoarece evidențiază atât avantajele, cât și limitările metodei.

\subsection{Completitudine}

MiniMax este un algoritm complet.

Dacă arborele de joc este finit și poate fi explorat integral, algoritmul va găsi întotdeauna o soluție.

Cu alte cuvinte, dacă există o mutare optimă, MiniMax o va determina.

\subsection{Optimalitate}

MiniMax este un algoritm optim.

Presupunând că:

\begin{itemize}
\item ambii jucători joacă perfect;
\item arborele este explorat complet;
\item valorile terminale sunt corecte;
\end{itemize}

algoritmul returnează strategia optimă pentru jucătorul MAX.

Rezultatul obținut reprezintă cea mai bună valoare garantată.

\subsection{Strategie conservatoare}

MiniMax este un algoritm pesimist.

Atunci când analizează o mutare, MAX presupune întotdeauna că adversarul va alege răspunsul cel mai nefavorabil pentru el.

Astfel, algoritmul nu caută cel mai bun rezultat posibil, ci cel mai bun rezultat garantat.

\subsection{Complexitatea în timp}

Fie:

\begin{itemize}
\item b = factorul de ramificare;
\item d = adâncimea arborelui.
\end{itemize}

Complexitatea temporală este:

\[
O(b^d)
\]

Numărul de stări analizate crește exponențial cu adâncimea.

Exemplu:

\[
b=10,\quad d=6
\]

Numărul aproximativ de noduri analizate este:

\[
10^6 = 1\,000\,000
\]

Acest cost devine rapid foarte mare pentru jocurile complexe.

\subsection{Complexitatea în spațiu}

MiniMax utilizează o explorare de tip Depth First Search (DFS).

Prin urmare, complexitatea spațială este:

\[
O(b \cdot d)
\]

Memoria necesară este semnificativ mai mică decât timpul de execuție.

\subsection{Avantaje}

Principalele avantaje sunt:

\begin{itemize}
\item garantează alegerea mutării optime;
\item are o fundamentare matematică riguroasă;
\item poate fi implementat relativ ușor;
\item reprezintă baza multor algoritmi moderni pentru jocuri.
\end{itemize}

\subsection{Dezavantaje}

Principalele dezavantaje sunt:

\begin{itemize}
\item cost computațional foarte ridicat;
\item explorarea completă a arborelui este adesea imposibilă;
\item timpul crește exponențial cu adâncimea;
\item devine impracticabil pentru jocuri cu factor mare de ramificare.
\end{itemize}

\subsection{Exemplu: Șah}

În șah, factorul mediu de ramificare este aproximativ:

\[
b \approx 35
\]

Pentru o adâncime de numai 8 mutări:

\[
35^8 \approx 2.25 \times 10^{12}
\]

Numărul de poziții este enorm.

Explorarea completă devine imposibilă chiar și pe calculatoare foarte puternice.

\subsection{Consecințe practice}

Din cauza complexității ridicate, MiniMax este rar utilizat singur.

În practică sunt introduse tehnici suplimentare:

\begin{itemize}
\item limitarea adâncimii de căutare;
\item funcții de evaluare;
\item Alpha-Beta Pruning.
\end{itemize}

Aceste tehnici reduc dramatic numărul de stări analizate.

\subsection{Rezumat}

\begin{center}
\begin{tabular}{|l|l|}
\hline
Proprietate & Valoare \\
\hline
Completitudine & Da \\
\hline
Optimalitate & Da \\
\hline
Complexitate timp & $O(b^d)$ \\
\hline
Complexitate spațiu & $O(b \cdot d)$ \\
\hline
Strategie & MAX vs MIN \\
\hline
\end{tabular}
\end{center}

MiniMax este un algoritm complet și optim, însă costul său exponențial reprezintă principala limitare și motivația introducerii tehnicii Alpha-Beta Pruning.







\enddocument{}